O câncer de próstata é uma das neoplasias malignas mais comuns em homens em todo o mundo, representando um desafio significativo para a saúde pública~\cite{ref-siegel2024}. O diagnóstico preciso e a graduação do câncer de próstata são cruciais para determinar estratégias de tratamento apropriadas e o prognóstico do paciente. O sistema de graduação da Sociedade Internacional de Patologia Urológica (ISUP), que classifica o câncer de próstata em graus 0--5 com base em padrões histológicos, serve como padrão-ouro para avaliação prognóstica~\cite{ref-epstein2016}.

A análise histopatológica tradicional depende do exame manual de imagens de lâminas completas (WSIs) por patologistas experientes. No entanto, esse processo apresenta vários desafios: (1) variabilidade inter e intra-observador nas decisões de graduação~\cite{ref-bulten2020}, (2) análise demorada de grandes amostras de tecido, (3) a necessidade de expertise altamente especializada, e (4) o volume crescente de biópsias que precisam ser avaliadas. Essas limitações motivaram o desenvolvimento de sistemas automatizados de patologia computacional baseados em aprendizado profundo.

Redes neurais convolucionais (CNNs) demonstraram sucesso notável na análise de imagens médicas, particularmente em histopatologia~\cite{ref-campanella2019}. Avanços recentes em arquiteturas de aprendizado profundo mostraram resultados promissores para a graduação automatizada de Gleason e classificação ISUP do câncer de próstata~\cite{ref-nagpal2019,ref-bulten2020}. Entre essas arquiteturas, o EfficientNet~\cite{ref-tan2019} emergiu como uma família de modelos particularmente eficaz, alcançando desempenho de estado da arte com significativamente menos parâmetros em comparação com arquiteturas tradicionais como ResNet~\cite{ref-he2016}.

O EfficientNet emprega um método de escalonamento composto que equilibra uniformemente a profundidade, largura e resolução da rede, resultando em eficiência e precisão aprimoradas. Essa característica o torna particularmente adequado para aplicações de imagens médicas onde os recursos computacionais podem ser limitados, mas alta precisão é necessária. A eficiência da arquitetura é alcançada através de blocos de convolução invertida em gargalo móvel (MBConv) com otimização squeeze-and-excitation~\cite{ref-hu2018}.

Neste trabalho, apresentamos um estudo abrangente sobre classificação automatizada de grau ISUP do câncer de próstata usando arquiteturas EfficientNet. Avaliamos múltiplas variantes do EfficientNet (B0, B3 e B7) no conjunto de dados do Desafio Prostate cANcer graDe Assessment (PANDA)~\cite{ref-panda2020}, que contém mais de 10.000 imagens de lâminas completas anotadas de múltiplas instituições. Nossa abordagem incorpora transfer learning a partir de pesos pré-treinados no ImageNet, estratégias de fine-tuning, customização de função de perda e métodos de ensemble para melhorar o desempenho de classificação.

Investigamos vários aspectos do pipeline de classificação: (1) o impacto de diferentes variantes do EfficientNet na precisão de classificação, (3) técnicas avançadas de treinamento incluindo remoção de ruídos baseada em entropia, média estocástica de pesos (SWA)~\cite{ref-izmailov2018}, perda focal, e (4) estratégias de ensemble combinando múltiplos modelos para alcançar predições robustas.

Nossas contribuições incluem: (1) uma avaliação sistemática de arquiteturas EfficientNet para graduação de câncer de próstata, (2) implementação e comparação de técnicas avançadas de treinamento, (3) desenvolvimento de métodos de ensemble que alcançam precisão de classificação aprimorada, (4) implementação da função de loss combinando perda ordinal e focal. Avaliamos nossos métodos usando kappa quadrático ponderado (QWK)~\cite{ref-cohen1968}, que é a métrica padrão para tarefas de classificação ordinal em patologia, juntamente com métricas de acurácia, F1-score, precisão e recall.
